\section{AIME I 2005}
        \begin{enumerate}
        	\item (\textit{AIME I 2005}) Six congruent circles form a ring with each circle externally tangent to two circles adjacent to it. All circles are internally tangent to a circle $C$ with radius 30. Let $K$ be the area of the region inside circle $C$ and outside of the six circles in the ring. Find $\lfloor K \rfloor.$


 
 \textbf{Answer: }942



\textbf{Solution}

Define the radii of the six congruent circles as $r$. If we draw all of the radii to the points of external tangency, we get a \href{https://artofproblemsolving.com/wiki/index.php?title=Regular\_polygon}{regular} \href{https://artofproblemsolving.com/wiki/index.php?title=Hexagon}{hexagon}. If we connect the \href{https://artofproblemsolving.com/wiki/index.php?title=Vertex}{vertices} of the hexagon to the \href{https://artofproblemsolving.com/wiki/index.php?title=Center}{center} of the circle $C$, we form several \href{https://artofproblemsolving.com/wiki/index.php?title=Equilateral\_triangle}{equilateral triangles}. The length of each side of the triangle is $2r$. Notice that the radius of circle $C$ is equal to the length of the side of the triangle plus $r$. Thus, the radius of $C$ has a length of $3r = 30$, and so $r = 10$. $K = 30^2\pi - 6(10^2\pi) = 300\pi$, so $\lfloor 300\pi \rfloor = \boxed{942}$.

 




	\item (\textit{AIME I 2005}) For each positive integer $k,$ let $S_k$ denote the increasing arithmetic sequence of integers whose first term is 1 and whose common difference is $k.$ For example, $S_3$ is the sequence $1,4,7,10,\ldots.$ For how many values of $k$ does $S_k$ contain the term 2005?


 
 \textbf{Answer: }012



\textbf{Solution}

Suppose that the $n$th term of the sequence $S_k$ is $2005$. Then $1+(n-1)k=2005$ so $k(n-1)=2004=2^2\cdot 3\cdot 167$.  The \href{https://artofproblemsolving.com/wiki/index.php?title=Ordered\_pair}{ordered pairs} $(k,n-1)$ of positive integers that satisfy the last equation are $(1,2004)$,$(2,1002)$, $(3,668)$, $(4,501)$, $(6,334)$, $(12,167)$, $(167,12)$,$(334,6)$, $(501,4)$, $(668,3)$, $(1002,2)$ and $(2004,1)$, and each of these gives a possible value of $k$. Thus the requested number of values is $12$, and the answer is $\boxed{012}$.

 Alternatively, notice that the formula for the number of \href{https://artofproblemsolving.com/wiki/index.php?title=Divisor}{divisors} states that there are $(2 + 1)(1 + 1)(1 + 1) = 12$ divisors of $2^2\cdot 3^1\cdot 167^1$.

 





\textbf{Solution 2}

Any term in the sequence $S_k$ can be written as 1+kx.  If this is to equal 2005, then the remainder when 2005 is divided by k is 1.

 Now all we have to do is find the numbers of factors of 2004.  There are $(2 + 1)(1 + 1)(1 + 1) = \boxed{012}$ divisors of $2^2\cdot 3^1\cdot 167^1$.

 Note that although the remainder when 2005 divided by 1 is not 1, it still works- $S_1$ would be the sequence of all positive integers, in which 2005 must appear.

 




	\item (\textit{AIME I 2005}) How many positive integers have exactly three proper divisors (positive integral divisors excluding itself), each of which is less than 50?


 
 \textbf{Answer: }109



\textbf{Solution (Basic Casework and Combinations)}

Suppose $n$ is such an \href{https://artofproblemsolving.com/wiki/index.php?title=Integer}{integer}. Because $n$ has $3$ proper divisors, it must have $4$ divisors,, so $n$ must be in the form $n=p\cdot q$ or $n=p^3$ for distinct \href{https://artofproblemsolving.com/wiki/index.php?title=Prime\_number}{prime numbers} $p$ and $q$. 

 In the first case, the three proper divisors of $n$ are $1$, $p$ and $q$.  Thus, we need to pick two prime numbers less than $50$. There are fifteen of these ($2, 3, 5, 7, 11, 13, 17, 19, 23, 29, 31, 37, 41, 43$ and $47$) so there are ${15 \choose 2} =105$ ways to choose a pair of primes from the list and thus $105$ numbers of the first type.

 In the second case, the three proper divisors of $n$ are 1, $p$ and $p^2$.  Thus we need to pick a prime number whose square is less than $50$.  There are four of these ($2, 3, 5,$ and $7$) and so four numbers of the second type. 

 Thus there are $105+4=\boxed{109}$ integers that meet the given conditions.

 




	\item (\textit{AIME I 2005}) The director of a marching band wishes to place the members into a formation that includes all of them and has no unfilled positions. If they are arranged in a square formation, there are 5 members left over. The director realizes that if he arranges the group in a formation with 7 more rows than columns, there are no members left over. Find the maximum number of members this band can have.


 
 \textbf{Answer: }294



\textbf{Solution 1}

If $n > 14$ then $n^2 + 6n + 14 < n^2 + 7n < n^2 + 8n + 21$ and so $(n + 3)^2 + 5 < n(n + 7) < (n + 4)^2 + 5$.  If $n$ is an \href{https://artofproblemsolving.com/wiki/index.php?title=Integer}{integer} there are no numbers which are 5 more than a \href{https://artofproblemsolving.com/wiki/index.php?title=Perfect\_square}{perfect square} strictly between $(n + 3)^2 + 5$ and $(n + 4)^2 + 5$.  Thus, if the number of columns is $n$, the number of students is $n(n + 7)$ which must be 5 more than a perfect square, so $n \leq 14$.  In fact, when $n = 14$ we have $n(n + 7) = 14\cdot 21 = 294 = 17^2 + 5$, so this number works and no larger number can.  Thus, the answer is $\boxed{294}$.

 





\textbf{Solution 2}

Define the number of rows/columns of the square formation as $s$, and the number of rows of the rectangular formation $r$ (so there are $r - 7$ columns). Thus, $s^2 + 5 = r(r-7) \Longrightarrow r^2 - 7r - s^2 - 5 = 0$. The \href{https://artofproblemsolving.com/wiki/index.php?title=Quadratic\_formula}{quadratic formula} yields $r = \frac{7 \pm \sqrt{49 - 4(1)(-s^2 - 5)}}{2} = \frac{7 \pm \sqrt{4s^2 + 69}}{2}$. $\sqrt{4s^2 + 69}$ must be an \href{https://artofproblemsolving.com/wiki/index.php?title=Integer}{integer}, say $x$. Then $4s^2 + 69 = x^2$ and $(x + 2s)(x - 2s) = 69$. The factors of $69$ are $(1,69), (3,23)$; $x$ is maximized for the first case. Thus, $x = \frac{69 + 1}{2} = 35$, and $r = \frac{7 \pm 35}{2} = 21, -14$. The latter obviously can be discarded, so there are $21$ rows and $21 - 7 = 14$ columns, making the answer $294$.

 





\textbf{Solution 3}

The number of members is $m^2+5=n(n+7)$ for some $n$ and $m$. Multiply both sides by $4$ and \href{https://artofproblemsolving.com/wiki/index.php?title=Completing\_the\_square}{complete the square} to get $4m^2+69=(2n+7)^2$. Thus, we have $69=((2n+7)+2m)((2n+7)-2m)$. Since we want to maximize $n$, set the first factor equal to $69$ and the second equal to $1$. Now we have the systems of equations 
 \[2n+2m+7 = 69\]
 \[2n-2m+7 = 1\]Adding these two equations gives us $4n+14=70$ which means $4n=56$ so $n=14$. Thus, the answer is $14\cdot21=294$. $\square$

 





\textbf{Solution 4}

Partially completing the square

 Geometrically: Split up the formation of $n + 7$ rows and $n$ columns into a square of $n$ rows and $n$ columns and a separate rectangle of the dimensions $7$ rows by $n$ columns. We want to take the rows from the rectangle and add them to the square to get another square and $5$ left over. If we attach exactly $2$ rows on the top and exactly $2$ rows on the side of the $n$ x $n$ square, then we have an $(n + 2)$ x $(n + 2)$ square that's missing a $2$ x $2$ corner. For the remaining $3n$ to fill this square plus the $5$ extra members, $n$ must be $3$. If we instead plaster exactly $3$ rows from the $7$ x $n$ formation to two adjacent sides of the $n$ x $n$ square, we have an $(n + 3)$ x $(n + 3)$ formation that's missing a $3$ x $3$ corner. For the remaining row of length $n$ to fill this plus five, $n = 14$. Plugging these in, we find $n = 14$ has a much higher count of members: $(n + 7)n; n = 14 --> 21(14) = 294$

 Algebraically: We have $n^2 + 7n = m$, where $m$ is the number of members in the band and $n$ is a positive integer. We partially complete the square for $n$ to get $n^2 + 7n = (n + 1)^2 + 5n - 1 = (n + 2)^2 + 3n - 4 = (n + 3)^2 + n - 9$Our goal is to get $n^2 + 7n = y^2 + 5$ because we want $m$ to be $5$ more than a perfect square. From the above, $5n - 1 = 5$ means $n$ isn't an integer, $3n - 4 = 5$ means that $n = 3$, and $n - 9 = 5$ means that $n = 14$. Out of these, $n = 14$ is associated with the highest number of members in the band, so $m = (14^2) + 7(14) = 294$

 


 





\textbf{Solution 5}

Let there be $m$ members and $n$ members for the square and $c$ for the number of columns of the other formation. We have $n^2 +5 = c(c+7) \implies n^2+5 = \left(c+\frac{7}{2}\right)^2 -\frac{49}{4} \implies n^2 - \left(c+\frac{7}{2}\right)^2 = -\frac{69}{4} \implies \left(n-c-\frac{7}{2}\right)\left(n + c +\frac{7}{2}\right) \implies (2n-2c-7)(2n+2n+7) = -69.$

 To maximize this we let $2n+2c+7 = 69$ and $2n-2c-7 = -1.$ Solving we find $n = 17$ so the desired number of members is $17^2 + 5 = \boxed{294}.$

 





\textbf{Solution 6 (NO ALGEBRA)}

Think of the process of moving people from the last column to new rows. Since there are less columns than rows, for each column removed, there are people discarded to the "extra" pile to be placed at the end. To maximize the number of "extra" people to fill in the last few rows. We remove 3 columns and add 4 rows. For the first new row, one more person will be discarded. For the second, 1 extra person are added since there is one more row now, and there are 2 less columns. Thus, there are 3 extra people discarded. Similarly, 5 extra people are discarded for the third column. Now there are $5+1+3+5=14$ people in the extra pile to put as the last row, so there are $14(14+7)=\boxed{294}$ people.

 





\textbf{Solution 7 (Bash-not advisable)}

Note: Only do this if you have a LOT of time (and you've memorized all your perfect squares up to 1000).

 We can see that the number of members in the band must be of the form $n(n + 7)$ for some positive integer $n$. When $n = 28$, this product is $980$, and since AIME answers are nonnegative integers less than $1000$, we don't have to check any higher $n$. Also, we know that this product must be 5 more than a perfect square, so we can make a table as shown and bash:
 \[\begin{tabular}{|c|c|c|} n & n(n+7) & 5 more than a perfect square?\\ \hline 1 & 8 & no\\ \hline 2 & 18 & no\\ \hline 3 & 30 & yes\\ \hline 4 & 44 & no\\ \hline 5 & 60 & no\\ \hline 6 & 78 & no\\ \hline 7 & 98 & no\\ \hline 8 & 120 & no\\ \hline 9 & 144 & no\\ \hline 10 & 170 & no\\ \hline 11 & 198 & no\\ \hline 12 & 228 & no\\ \hline 13 & 260 & no\\ \hline 14 & 294 & yes\\ \hline 15 & 330 & no\\ \hline 16 & 368 & no\\ \hline 17 & 408 & no\\ \hline 18 & 450 & no\\ \hline 19 & 494 & no\\ \hline 20 & 540 & no\\ \hline 21 & 588 & no\\ \hline 22 & 638 & no\\ \hline 23 & 690 & no\\ \hline 24 & 744 & no\\ \hline 25 & 800 & no\\ \hline 26 & 858 & no\\ \hline 27 & 918 & no\\ \hline 28 & 980 & no\\ \hline \end{tabular}\]Thus, we can see that our largest valid $n(n+7)$ is $\boxed{294}$.

 




	\item (\textit{AIME I 2005}) Robert has 4 indistinguishable gold coins and 4 indistinguishable silver coins. Each coin has an engraving of one face on one side, but not on the other. He wants to stack the eight coins on a table into a single stack so that no two adjacent coins are face to face. Find the number of possible distinguishable arrangements of the 8 coins.


 
 \textbf{Answer: }630



\textbf{Solution 1}

There are two separate parts to this problem: one is the color (gold vs silver), and the other is the orientation.

 There are ${8\choose4} = 70$ ways to position the gold coins in the stack of 8 coins, which determines the positions of the silver coins.

 Create a string of letters H and T to denote the orientation of the top of the coin. To avoid making two faces touch, we cannot have the arrangement HT. Thus, all possible configurations must be a string of tails followed by a string of heads, since after the first H no more tails can appear. The first H can occur in a maximum of eight times different positions, and then there is also the possibility that it doesn’t occur at all, for $9$ total configurations. Thus, the answer is $70 \cdot 9 = \boxed{630}$.

 





\textbf{Solution 2}

We can imagine the $8$ coins as a string of $0\text{'s}$ and $1\text{'s}$. Because no $2$ adjacent coins can have $2$ faces touching, subsequent to changing from $0$ to $1$, the numbers following $1$ must be $1\text{'s}$; therefore, the number of possible permutations if all the coins are indistinguishable is $9$ (there are $8$ possible places to change from $0$ to $1$ and there is the possibility that there no change occurs). There are $\binom 8 4$ possibilities of what coins are gold and what coins are silver, so the solution is $\boxed{9\cdot \binom 8 4=630}$.

 





\textbf{Solution 3}

First, we can break this problem up into two parts: the amount of ways to order the coins based on color, and which side is facing up for each coin (assuming they are indistinguishable). In the end, we multiply these values together. 

 First, there are obviously $\binom{8}{4}$ ways to order the coins based on color. 

 Next, we set up a recurrence. Let $a_n$ be the number of ways $n$ indistinguishable coins to be stacked such that no heads are facing each other. 

 Consider the top coin. If it is facing up, then there are $a_{n-1}$ ways for the rest of the coins below it to be arranged. If it is facing down however, there will only be one way to arrange the coins. We can the recurrence: 
 \[a_n=a_{n-1}+1.\] Note $a_1=2$, so $a_8 = 9$.

 Finally, to get the result, we do $70\times 9$ to get $\boxed{630}$.

 \textbf{Intuition} Really try not to overthink especially in combinatorics. In this case, we just have to notice that after we place a coin facing up ANYWHERE, the rest of the coins to the right of the row of coins MUST be facing up. (Or else we have a coin facing up followed by a coin facing down, and that is not cool). So for example, suppose we started with a coin facing up. Then the next coin must also be facing up. Then so must the next coin. And so on. Now we know what happens if it starts with a coin facing up. Hmmmmmmm. What if it starts with a coin facing down? 

 Well, start with a coin facing down. Now, does the next coin HAVE to have a coin facing down, or can it have a coin facing up? It could have both! But wait! What happens if the next coin is heads? Well we saw already that now the next coin must be heads, or else we have a up coin followed by a down coin, which once again, sucks. So now everything to the right of the up coin must be facing up. 

 So we can have the following configurations: $9U, 1D+8U, 2D+7U, 3D+6U, ..., 9D+0U$. These are the total number of configurations, which is $9$. (As a sidenote, we really count the number of solutions to $a+b=8$, which is $8+1=9$). So we must now put the 4 gold coins as a part of this sequence of ups and downs. This can be done in $8\choose 4$ ways, and the rest of the silver coins can be put in $8-4$choose $4$ ways. (We subtract the $4$ from the $8$ for the silver coin placement, because we already used up $4$ coins, namely the gold ones). So we have the answer as $9 \binom {8}{4}=630$.

 




	\item (\textit{AIME I 2005}) Let $P$ be the product of the nonreal roots of $x^4-4x^3+6x^2-4x=2005.$ Find $\lfloor P\rfloor.$


 
 \textbf{Answer: }045



\textbf{Solution 1}

The left-hand side of that \href{https://artofproblemsolving.com/wiki/index.php?title=Equation}{equation} is nearly equal to $(x - 1)^4$.  Thus, we add 1 to each side in order to complete the fourth power and get$(x - 1)^4 = 2006$.

 Let $r = \sqrt[4]{2006}$ be the positive \href{https://artofproblemsolving.com/wiki/index.php?title=Real}{real} fourth root of 2006.  Then the roots of the above equation are $x = 1 + i^n r$ for $n = 0, 1, 2, 3$.  The two non-real members of this set are $1 + ir$ and $1 - ir$.  Their product is $P = 1 + r^2 = 1 + \sqrt{2006}$.  $44^2 = 1936 < 2006 < 2025 = 45^2$ so $\lfloor P \rfloor = 1 + 44 = \boxed{045}$.

 





\textbf{Solution 2}

Starting like before,$(x-1)^4= 2006$This time we apply differences of squares.$(x-1)^4-2006=0$ so$((x-1)^2+\sqrt{2006})((x-1)^2 -\sqrt{2006})=0$If you think of each part of the product as a quadratic, then $((x-1)^2+\sqrt{2006})$ is bound to hold the two  non-real roots since the other definitely crosses the x-axis twice since it is just $x^2$ translated down and right.Therefore $P$ is the product of the roots of $((x-1)^2+\sqrt{2006})$ or $P=1+\sqrt{2006}$ so

 $\lfloor P \rfloor = 1 + 44 = \boxed{045}$.

 





\textbf{Solution 3}

If we don't see the fourth power, we can always factor the LHS to try to create a quadratic substitution. Checking, we find that $x=0$ and $x=2$ are both roots. Synthetic division gives$(x^2-2x)(x^2-2x+2)=2005$. We now have our quadratic substitution of$y=x^2-2x+1=(x-1)^2$, giving us$(y-1)(y+1)=2005$. From here we proceed as in Solution 1 to get $\boxed{045}$.

 





\textbf{Solution 4}

Realizing that if we add 1 to both sides we get $x^4-4x^3+6x^2-4x+1=2006$ which can be factored as $(x-1)^4=2006$. Then we can substitute $(x-1)$ with $y$ which leaves us with $y^4=2006$. Now subtracting 2006 from both sides we get some difference of squares $y^4-2006=0 \rightarrow (y-\sqrt[4]{2006})(y+\sqrt[4]{2006})(y^2+\sqrt{2006})=0$. The question asks for the product of the complex roots so we only care about the last factor which is equal to zero. From there we can solve $y^2+\sqrt{2006}=0$, we can substitute $(x-1)$ for $y$ giving us  $(x-1)^2+\sqrt{2006}=0$, expanding this we get $x^2-2x+1+\sqrt{2006}=0$. We know that the product of a quadratics roots is $\frac{c}{a}$ which leaves us with $\frac{1+\sqrt{2006}}{1}=1+\sqrt{2006}\approx\boxed{045}$.

 





\textbf{Solution 5}

As in solution 1, we find that $(x-1)^4 = 2006$. Now $x-1=\pm \sqrt[4]{2006}$ so $x_1 = 1+\sqrt[4]{2006}$ and $x_2 = 1-\sqrt[4]{2006}$ are the real roots of the equation. Multiplying, we get $x_1 x_2 = 1 - \sqrt{2006}$. Now transforming the original function and using Vieta's formula, $x^4-4x^3+6x^2-4x-2005=0$ so $x_1 x_2 x_3 x_4 = \frac{-2005}{1} = -2005$. We find that the product of the nonreal roots is $x_3 x_4 = \frac{-2005}{1-\sqrt{2006}} \approx 45.8$ and we get $\boxed{045}$.

 


 Note: $\frac{2005}{\sqrt[4]{2006}-1}=\frac{2005(1+\sqrt[4]{2006})}{2005} = 1+\sqrt[4]{2006}.$

 





\textbf{Solution 6 (De Moivre's Theorem)}

As all the other solutions, we find that $(x-1)^4 = 2006$. Thus $x=\sqrt[4]{2006}+1$. Thus $x= \sqrt[4]{2006}(\cos(\frac{2\pi(k)}{4}+i\sin(\frac{2\pi(k)}{4}))+1$ when $k=0,1,2,3$. The complex values of $x$ are the ones where $i\sin(\frac{2\pi(k)}{4})$ does not equal 0. These complex roots are $1+\sqrt[4]{2006}(i)$ and $1-\sqrt[4]{2006}(i)$. The product of these two nonreal roots is ($1+\sqrt[4]{2006}(i)$)($1-\sqrt[4]{2006}(i)$) which is equal to $1+\sqrt {2006}$. The floor of that value is $\boxed{045}$.

 





\textbf{Solution 7 (The Floor Divisor Method)}

Continue Solution 1 until we obtain the equation \( (x-1)^4 = 2006 \). Then, we see that \( x = 1 \pm \sqrt[4]{2006} \). While one could see the possible roots are \( 1 \pm \sqrt[4]{2006} \), and \( 1 \pm \sqrt[4]{2006}i \), and obtain 45, one could have also seen the product of all roots is \( -2005 \) through Vieta's. Then, the product of the two remaining real roots is \( 1 - \sqrt{2006} \), and thus the wanted value is \( \left\lfloor \frac{-2005}{1 - \sqrt{2006}} \right\rfloor \) or \( \left\lfloor \frac{2005}{\sqrt{2006} - 1} \right\rfloor \), which through the divisor method is \( \left\lfloor \frac{2005}{\lfloor \sqrt{2006} - 1 \rfloor} \right\rfloor +1 \), or \( \left\lfloor \frac{2005}{43} \right\rfloor +1 \), whose floor is $\boxed{045}$. 

 The formal Floor Divisor Method states that for all $a$ and $b$, where $A = \lfloor a \rfloor, \quad B = \lfloor b \rfloor$, the floor of $\frac{a}{b}$ is 
 \[\lfloor a/b \rfloor = \left\lfloor \frac{A}{B} - \frac{\mathbf{1}_{(a \bmod b) < (A \bmod B)}}{1} \right\rfloor\]

 More informally, without the bolded 1 (conditional 1), 

 For all $A = \lfloor a \rfloor, \quad B = \lfloor b \rfloor$, \begin{array}{c} \lfloor a/b \rfloor =\begin{cases} \lfloor A/B \rfloor & \text{if } a \ge \lfloor A/B \rfloor \cdot b \\[2mm]\lfloor A/B \rfloor - 1 & \text{if } a < \lfloor A/B \rfloor \cdot b\end{cases} \end{array}

 As a problem to the reader, try to figure out why I simplified this formula above into $\left\lfloor \frac{a}{b} \right\rfloor +1$ for this problem, and when should and should not simplify to this form.

 




	\item (\textit{AIME I 2005}) In quadrilateral $ABCD, BC=8, CD=12, AD=10,$ and $m\angle A= m\angle B = 60^\circ.$ Given that $AB = p + \sqrt{q},$ where $p$ and $q$ are positive integers, find $p+q.$


 
 \textbf{Answer: }150



\textbf{Solution 1}


                    \begin{figure}[H]
                        \centering
                        \includegraphics[width=0.5\textwidth]{Solution/AIME_I_2005_7_Solution_Figure_1.png}
                    \end{figure}
                    

 Draw line segment $DE$ such that line $DE$ is concurrent with line $BC$. Then, $ABED$ is an isosceles trapezoid so $AD=BE=10$, and $BC=8$ and $EC=2$. We are given that $DC=12$. Since $\angle CED = 120^{\circ}$, using Law of Cosines on $\bigtriangleup CED$ gives 
 \[12^2=DE^2+4-2(2)(DE)(\cos 120^{\circ})\] which gives 
 \[144-4=DE^2+2DE\]. Adding $1$ to both sides gives $141=(DE+1)^2$, so $DE=\sqrt{141}-1$. $\bigtriangleup DAP$ and $\bigtriangleup EBQ$ are both $30-60-90$, so $AP=5$ and $BQ=5$. $PQ=DE$, and therefore $AB=AP+PQ+BQ=5+\sqrt{141}-1+5=9+\sqrt{141} \rightarrow (p,q)=(9,141) \rightarrow \boxed{150}$.

 





\textbf{Solution 2}


                    \begin{figure}[H]
                        \centering
                        \includegraphics[width=0.5\textwidth]{Solution/AIME_I_2005_7_Solution_Figure_2.png}
                    \end{figure}
                    Draw the \href{https://artofproblemsolving.com/wiki/index.php?title=Perpendicular}{perpendiculars} from $C$ and $D$ to $AB$, labeling the intersection points as $E$ and $F$. This forms 2 $30-60-90$ \href{https://artofproblemsolving.com/wiki/index.php?title=Right\_triangle}{right triangles}, so $AE = 5$ and $BF = 4$. Also, if we draw the horizontal line extending from $C$ to a point $G$ on the line $DE$, we find another right triangle $\triangle DGC$. $DG = DE - CF = 5\sqrt{3} - 4\sqrt{3} = \sqrt{3}$. The \href{https://artofproblemsolving.com/wiki/index.php?title=Pythagorean\_Theorem}{Pythagorean Theorem} yields that $GC^2 = 12^2 - \sqrt{3}^2 = 141$, so $EF = GC = \sqrt{141}$. Therefore, $AB = 5 + 4 + \sqrt{141} = 9 + \sqrt{141}$, and $p + q = \boxed{150}$.

 





\textbf{Solution 3}


                    \begin{figure}[H]
                        \centering
                        \includegraphics[width=0.5\textwidth]{Solution/AIME_I_2005_7_Solution_Figure_3.png}
                    \end{figure}
                    Extend $AD$ and $BC$ to an intersection at point $E$. We get an \href{https://artofproblemsolving.com/wiki/index.php?title=Equilateral\_triangle}{equilateral triangle} $ABE$. We denote the length of a side of $\triangle ABE$ as $s$ and solve for it using the \href{https://artofproblemsolving.com/wiki/index.php?title=Law\_of\_Cosines}{Law of Cosines}: 
 \[12^2 = (s - 10)^2 + (s - 8)^2 - 2(s - 10)(s - 8)\cos{60}\]
 \[144 = 2s^2 - 36s + 164 - (s^2 - 18s + 80)\] This simplifies to $s^2 - 18s - 60=0$; the \href{https://artofproblemsolving.com/wiki/index.php?title=Quadratic\_formula}{quadratic formula} yields the (discard the negative result) same result of $9 + \sqrt{141}$.

 





\textbf{Solution 4}

Extend $BC$ and $AD$ to meet at point $E$, forming an equilateral triangle $\triangle ABE$. Draw a line from $C$ parallel to $AB$ so that it intersects $AD$ at point $F$. Then, apply \href{https://artofproblemsolving.com/wiki/index.php?title=Stewart\%27s\_Theorem}{Stewart's Theorem} on $\triangle CFE$. Let $CE=x$. 
 \[2x(x-2) + 12^2x = 2x^2 + x^2(x-2)\] 
 \[x^3 - 2x^2 - 140x = 0\] By the quadratic formula (discarding the negative result), $x = 1 + \sqrt{141}$, giving $AB = 9 + \sqrt{141}$ for a final answer of $p+q=150$.

 





\textbf{Solution 5}

Draw a line from point $C$ to a new point $E$ on $AD$ parallel to $AB$. Draw a line from point $E$ to a new point $F$ on $AB$ parallel to $CD$. This creates parallelogram $CEFB$. 

 \begin{verbatim}
The reasoning for this is to connect the two angles that are congruent. 
\end{verbatim}$\angle CEF = 60^{\circ}$ and triangle $AEF$ is an equilateral triangle with side length 8. Thus, $AF = 8$. 

 $ED = 10 - 8 = 2$. By the Law of Cosines, $CD^2 = ED^2 + EC^2 - 2 \cdot {ED} \cdot {EC} \cdot \cos{\angle DEC}$. Thus, $144 = 4 + c^2 - 2c$ and $c = 1 + \sqrt{141}$. $EC = FB$. 

 $AB = AF + FB = 8 + 1 + \sqrt{141} = 9 + \sqrt{141}$ so the answer is $p+q=150$.

 -unhappyfarmer

 




	\item (\textit{AIME I 2005}) The equation $2^{333x-2} + 2^{111x+2} = 2^{222x+1} + 1$ has three real roots. Given that their sum is $\frac mn$ where $m$ and $n$ are relatively prime positive integers, find $m+n.$


 
 \textbf{Answer: }113



\textbf{Solution}

Let $y = 2^{111x}$.  Then our equation reads $\frac{1}{4}y^3 + 4y = 2y^2 + 1$ or $y^3 - 8y^2 + 16y - 4 = 0$.  Thus, if this equation has roots $r_1, r_2$ and $r_3$, by \href{https://artofproblemsolving.com/wiki/index.php?title=Vieta\%27s\_formulas}{Vieta's formulas} we have $r_1\cdot r_2\cdot r_3 = 4$.  Let the corresponding values of $x$ be $x_1, x_2$ and $x_3$.  Then the previous statement says that $2^{111\cdot(x_1 + x_2 + x_3)} = 4$ so taking a \href{https://artofproblemsolving.com/wiki/index.php?title=Logarithm}{logarithm} of that gives $111(x_1 + x_2 + x_3) = 2$ and $x_1 + x_2 + x_3 = \frac{2}{111}$.  Thus the answer is $111 + 2 = \boxed{113}$.

 




	\item (\textit{AIME I 2005}) Twenty-seven unit cubes are painted orange on a set of four faces so that the two unpainted faces share an edge. The 27 cubes are then randomly arranged to form a $3\times 3 \times 3$ cube. Given that the probability that the entire surface of the larger cube is orange is $\frac{p^a}{q^br^c},$ where $p,q,$ and $r$ are distinct primes and $a,b,$ and $c$ are positive integers, find $a+b+c+p+q+r.$


 
 \textbf{Answer: }074



\textbf{Solution}


                    \begin{figure}[H]
                        \centering
                        \includegraphics[width=0.5\textwidth]{Solution/AIME_I_2005_9_Solution_Figure_1.png}
                    \end{figure}
                    

 We can consider the orientation of each of the individual cubes independently.The unit cube at the center of our large cube has no exterior faces, so all of its orientations work.

 For the six unit cubes and the centers of the faces of the large cube, we need that they show an orange face.  This happens in $\frac{4}{6} = \frac{2}{3}$ of all orientations, so from these cubes we gain a factor of $\left(\frac{2}{3}\right)^6$.

 The twelve unit cubes along the edges of the large cube have two faces showing, and these two faces are joined along an edge.  Thus, we need to know the number of such pairs that are both painted orange.  We have a pair for each edge, and 7 edges border one of the unpainted faces while only 5 border two painted faces.  Thus, the probability that two orange faces show for one of these cubes is $\frac{5}{12}$, so from all of these cubes we gain a factor of $\left(\frac{5}{12}\right)^{12} = \frac{5^{12}}{2^{24}3^{12}}$.

 Finally, we need to orient the eight corner cubes.  Each such cube has 3 faces showing, and these three faces share a common vertex.  Thus, we need to know the number of vertices for which all three adjacent faces are painted orange.  There are six vertices which are a vertex of one of the unpainted faces and two vertices which have our desired property, so each corner cube contributes a probability of $\frac{2}{8} = \frac{1}{4}$ and all the corner cubes together contribute a probability of $\left(\frac{1}{4}\right)^8 = \frac{1}{2^{16}}$

 Since these probabilities are independent, the overall probability is just their product, $\frac{2^6}{3^6} \cdot \frac{5^{12}}{2^{24}3^{12}} \cdot \frac{1}{2^{16}} = \frac{5^{12}}{2^{34}\cdot 3^{18}}$ and so the answer is $2 + 3 + 5 + 12 + 34 + 18 = \boxed{074}$.

 





\textbf{Solution 2}

Rather than worry about the actual painted faces, consider the position of the shared edge of the non-colored faces. The six centers of the faces can be calculated in the same manner as last time, getting $\left(\frac{2}{3}\right)^6$. 

 The twelve edge cubes show two faces. We want the shared edge of the non-colored faces to not be one of the edges that are exposed on the outside. There are 7 such edges that can be seen when viewed from the outside, so the probability they cannot be seen are $\left(\frac{5}{12}\right)^{12}$.

 The eight corner cubes each have 9 visible edges from the outside, so the probability that the edge of the non-colored faces is not one of those is $\left(\frac{3}{12}\right)^8 = \left(\frac{1}{4}\right)^8$. Multiplying the probabilities together, we get the same exact solution.  

 As a side note, notice that the placement of the two unpainted faces is in fact of vital importance: if they were on opposite faces, the answer would be 0 because any placement of such a cube in the corner of the large cube would show one unpainted face.

 




	\item (\textit{AIME I 2005}) Triangle $ABC$ lies in the Cartesian Plane and has an area of 70. The coordinates of $B$ and $C$ are $(12,19)$ and $(23,20),$ respectively, and the coordinates of $A$ are $(p,q).$ The line containing the median to side $BC$ has slope $-5.$ Find the largest possible value of $p+q.$


 
 \textbf{Answer: }047



\textbf{Solution 1}

The \href{https://artofproblemsolving.com/wiki/index.php?title=Midpoint}{midpoint} $M$ of \href{https://artofproblemsolving.com/wiki/index.php?title=Line\_segment}{line segment} $\overline{BC}$ is $\left(\frac{35}{2}, \frac{39}{2}\right)$. The equation of the median can be found by $-5 = \frac{q - \frac{39}{2}}{p - \frac{35}{2}}$. Cross multiply and simplify to yield that $-5p + \frac{35 \cdot 5}{2} = q - \frac{39}{2}$, so $q = -5p + 107$.

 Use \href{https://artofproblemsolving.com/wiki/index.php?title=Determinant}{determinants} to find that the \href{https://artofproblemsolving.com/wiki/index.php?title=Area}{area} of $\triangle ABC$ is $\frac{1}{2} \begin{vmatrix}p & 12 & 23 \\  q & 19 & 20 \\ 1 & 1 & 1\end{vmatrix} = 70$ (note that there is a missing \href{https://artofproblemsolving.com/wiki/index.php?title=Absolute\_value}{absolute value}; we will assume that the other solution for the triangle will give a smaller value of $p+q$, which is provable by following these steps over again) (alternatively, we could use the \href{https://artofproblemsolving.com/wiki/index.php?title=Shoelace\_Theorem}{Shoelace Theorem}). We can calculate this determinant to become $140 = \begin{vmatrix} 12 & 23 \\ 19 & 20 \end{vmatrix} - \begin{vmatrix} p & q \\ 23 & 20 \end{vmatrix} + \begin{vmatrix} p & q \\ 12 & 19 \end{vmatrix}$ $\Longrightarrow 140 = 240 - 437 - 20p + 23q + 19p - 12q$ $= -197 - p + 11q$. Thus, $q = \frac{1}{11}p + \frac{337}{11}$.

 Setting this equation equal to the equation of the median, we get that $\frac{1}{11}p + \frac{337}{11} = -5p + 107$, so $\frac{56}{11}p = \frac{107 \cdot 11 - 337}{11}$. Solving produces that $p = 15$. \href{https://artofproblemsolving.com/wiki/index.php?title=Substitution}{Substituting} backwards yields that $q = 32$; the solution is $p + q = \boxed{047}$.

 





\textbf{Solution 2}

Using the equation of the median from above, we can write the \href{https://artofproblemsolving.com/wiki/index.php?title=Coordinate}{coordinates} of $A$ as $(p,\ -5p + 107)$. The equation of $\overline{BC}$ is $\frac{20 - 19}{23 - 12} = \frac{y - 19}{x - 12}$, so $x - 12 = 11y - 209$. In \href{https://artofproblemsolving.com/wiki/index.php?title=General\_form&action=edit&redlink=1}{general form}, the line is $x - 11y + 197 = 0$. Use the equation for the distance between a line and point to find the distance between $A$ and $BC$ (which is the height of $\triangle ABC$): $\frac{|1(p) - 11(-5p + 107) + 197|}{1^2 + 11^2} = \frac{|56p - 980|}{\sqrt{122}}$. Now we need the length of $BC$, which is $\sqrt{(23 - 12)^2 + (20 - 19)^2} = \sqrt{122}$. The area of $\triangle ABC$ is $70 = \frac{1}{2}bh = \frac{1}{2}\left(\frac{|56p - 980|}{\sqrt{122}}\right) \cdot \sqrt{122}$. Thus, $|28p - 490| = 70$, and $p = 15,\ 20$. We are looking for $p + q = -4p + 107 = 47,\ 27$. The maximum possible value of $p + q = \fbox{047}$.

 





\textbf{Solution 3}

Again, the \href{https://artofproblemsolving.com/wiki/index.php?title=Midpoint}{midpoint} $M$ of \href{https://artofproblemsolving.com/wiki/index.php?title=Line\_segment}{line segment} $\overline{BC}$ is at $\left(\frac{35}{2}, \frac{39}{2}\right)$. Let $A'$ be the point $(17, 22)$, which lies along the line through $M$ of slope $-5$.  The area of triangle $A'BC$ can be computed in a number of ways (one possibility: extend $A'B$ until it hits the line $y = 19$, and subtract one triangle from another), and each such calculation gives an area of 14.  This is $\frac{1}{5}$ of our needed area, so we simply need the point $A$ to be 5 times as far from $M$ as $A'$ is.  Thus $A = \left(\frac{35}{2}, \frac{39}{2}\right) \pm 5\left(-\frac{1}{2}, \frac{5}{2}\right)$, and the sum of coordinates will be larger if we take the positive value, so $A = \left(\frac{35}{2} - \frac{5}2, \frac{39}{2} + \frac{25}{2}\right)$ and the answer is $\frac{35}{2} - \frac{5}2 + \frac{39}{2} + \frac{25}{2} = \fbox{047}$.

 





\textbf{Solution 4}

Plug points into the \href{https://artofproblemsolving.com/wiki/index.php?title=Shoelace\_Theorem}{Shoelace Theorem}. This will provide you with the equation $|-p - 197 + 11q| = 140$. The find the midpoint of the line $BC$ which is $(17.5,19.5)$. Now using this post and the given slope of the median, $-5$, using basic algebra we can find the equation of the median which is $q = -5p + 107$. Now that we have been given $q$ in terms of $p$ plug this equation back into $|-p - 197 + 11q| = 140$. The result is the equation $|980 - 56p| = 140$. Solve this equation for two possible answers $p = 15, 20$. Plugging into $q = -5p + 107$ these inputs produce $q$ values $32$ and $7$. Obviously $15 + 32$ is the greater sum so the answer is $47$ and we are done.

 




	\item (\textit{AIME I 2005}) A semicircle with diameter $d$ is contained in a square whose sides have length 8. Given the maximum value of $d$ is $m - \sqrt{n},$ find $m+n.$


 
 \textbf{Answer: }544



\textbf{Solution}

We note that aligning the base of the semicircle with a side of the square is certainly non-optimal.  If the semicircle is tangent to only one side of the square, we will have "wiggle-room" to increase its size.  Once it is tangent to two adjacent sides of the square, we will maximize its size when it touches both other sides of the square.  This can happen only when it is arranged so that the center of the semicircle lies on one diagonal of the square.

 
                    \begin{figure}[H]
                        \centering
                        \includegraphics[width=0.5\textwidth]{Solution/AIME_I_2005_11_Solution_Figure_1.png}
                    \end{figure}
                    

 Now, let the square be $ABCD$, and let $E \in AB$ and $F \in DA$ be the points at which the "corners" of the semicircle touch the square. Let $O$ be the center of the semicircle.

 





\textbf{Solution 1}

We can just look at a quarter circle inscribed in a $45-45-90$ right triangle. We can then extend a radius, $r$ to one of the sides creating an $r,r, r\sqrt{2}$ right triangle. This means that we have $r + r\sqrt{2} = 8\sqrt{2}$ so $r = \frac{8\sqrt{2}}{1+\sqrt{2}} = 16 - 8\sqrt{2}$. Then the diameter is $32 - \sqrt{512}$ giving us $32 + 512 = \boxed{544}$

 





\textbf{Solution 2}

Define the radius of the semicircle as $r$. Draw the \href{https://artofproblemsolving.com/wiki/index.php?title=Perpendicular}{perpendicular} from $O$ to $AB$, which forms a $45-45-90$ triangle. The length of the perpendicular is $\frac{r}{\sqrt{2}}$. Note also that $AD$ is equal to the length of that perpendicular plus the radius to the point of tangency on $CD$. Thus, $r + \frac{r}{\sqrt{2}} = 8$, and $r = \frac{8\sqrt{2}}{\sqrt{2} + 1} \cdot (\frac{\sqrt{2} - 1}{\sqrt{2} - 1}) = 16 - 8\sqrt{2}$. The diameter is then $2r = 32 - \sqrt{512}$, and the solution is $32 + 512 = 544$.

 





\textbf{Solution 3}

By the comments above, $AE = AF = a$.  By the \href{https://artofproblemsolving.com/wiki/index.php?title=Pythagorean\_Theorem}{Pythagorean Theorem}, $d^2 = 2a^2$.

 Now, if we draw a line through the center, $O$, of the semicircle and its point of tangency with $BC$, we see that this line is perpendicular to $BC$ and so parallel to $AB$.  Thus, by \href{https://artofproblemsolving.com/wiki/index.php?title=Triangle\_similarity&action=edit&redlink=1}{triangle similarity} it cuts $AF$ in half, and so by symmetry the distance from $O$ to $AD$ is $\frac{a}{2}$ and so the distance from $O$ to $BC$ is $8 - \frac a2$.  But this latter quantity is also the radius of the semicircle, so $d = 16 - a$.

 Our two previous paragraphs give $2a^2 = (16 - a)^2$ so $a^2 + 32a - 256 = 0$ and $a = 16\sqrt{2} - 16$ (where we discard the negative root of that quadratic) and so $d = a\sqrt{2} = 32 - 16\sqrt{2} = 32 - \sqrt{512}$, so the answer is $32 + 512 = 544$.

 





\textbf{Solution 4}

We proceed by finding the area of the square in 2 different ways. The square is obviously 8*8=64, but we can also find the area in terms of d. From the center of the circle, draw radii that hit the points where the square is tangent to the semicircle. Then the square's area is the area of the small square +2* the area of the trapezoids on the corners+ the area of an isoceles triangle. Adding these all up gives $64=\frac{d^2}{4}+\frac{d^2}{4}+(8-\frac{d\sqrt{2}}{2}+\frac{d}{2})(8-\frac{d}{2})$Simplifying gives $d-16\sqrt{2}+d\sqrt{2}=0$. Solving gives $d=32-16\sqrt{2}=32-\sqrt{512}$ so the answer is $32+512= \boxed{544}$.

 





\textbf{Solution 5}

It is easy after getting the image, after drawing labeling the lengths of those segments, assume the radius is $x$, we can see $x=\sqrt{2}(8-x)$ and we get $2x=32-\sqrt{512}$ and we have the answer $\boxed{544}$ ~bluesoul

 




	\item (\textit{AIME I 2005}) For positive integers $n,$ let $\tau (n)$ denote the number of positive integer divisors of $n,$ including 1 and $n.$ For example, $\tau (1)=1$ and $\tau(6) =4.$ Define $S(n)$ by $S(n)=\tau(1)+ \tau(2) + \cdots + \tau(n).$ Let $a$ denote the number of positive integers $n \leq 2005$ with $S(n)$ odd, and let $b$ denote the number of positive integers $n \leq 2005$ with $S(n)$ even. Find $|a-b|.$


 
 \textbf{Answer: }025



\textbf{Solution 1}

It is well-known that $\tau(n)$ is odd if and only if $n$ is a \href{https://artofproblemsolving.com/wiki/index.php?title=Perfect\_square}{perfect square}.  (Otherwise, we can group \href{https://artofproblemsolving.com/wiki/index.php?title=Divisor}{divisors} into pairs whose product is $n$.)  Thus, $S(n)$ is odd if and only if there are an odd number of perfect squares less than $n$.  So $S(1), S(2)$ and $S(3)$ are odd, while $S(4), S(5), \ldots, S(8)$ are even, and $S(9), \ldots, S(15)$ are odd, and so on. 

 So, for a given $n$, if we choose the positive integer $m$ such that $m^2 \leq n < (m + 1)^2$ we see that $S(n)$ has the same parity as $m$.  

 It follows that the numbers between $1^2$ and $2^2$, between $3^2$ and $4^2$, and so on, all the way up to the numbers between $43^2$ and $44^2 = 1936$ have $S(n)$ odd. These are the only such numbers less than $2005$ (because $45^2 = 2025 > 2005$).

 





\textbf{Solution 2}

Notice that the difference between consecutive squares are consecutively increasing odd numbers. Thus, there are $3$ numbers between $1$ (inclusive) and $4$ (exclusive), $5$ numbers between $4$ and $9$, and so on. The number of numbers from $n^2$ to $(n + 1)^2$ is $(n + 1 - n)(n + 1 + n) = 2n + 1$. Whenever the lowest square beneath a number is odd, the parity will be odd, and the same for even. Thus, $a = [2(1) + 1] + [2(3) + 1] \ldots [2(43) + 1] = 3 + 7 + 11 \ldots 87$. $b = [2(2) + 1] + [2(4) + 1] \ldots [2(42) + 1] + 70 = 5 + 9 \ldots 85 + 70$, the $70$ accounting for the difference between $2005$ and $44^2 = 1936$, inclusive. Notice that if we align the two and subtract, we get that each difference is equal to $2$. Thus, the solution is $|a - b| = |b - a| = |2 \cdot 21 + 70 - 87| = \boxed{025}$.

 





\textbf{Solution 3}

Similarly, $b = (3^2 - 2^2) + (5^2 - 4^2) + \ldots + (45^2 - 44^2) - 19$, where the $-19$ accounts for those numbers between $2005$ and $2024$.

 Thus $a = (2^2 - 1^2) + (4^2 - 3^2) + \ldots + (44^2 - 43^2)$.

 Then, $|a - b| = |2(2^2 + 4^2 + \ldots + 44^2) - 2(1^2 + 3^2 + 5^2 + \ldots 43^2) + 1^2 - 45^2 + 19|$.We can apply the formula $1^2 + 2^2 + \ldots + n^2 = \frac{n(n + 1)(2n + 1)}{6}$.  From this formula, it follows that $2^2 + 4^2 + \ldots + (2n)^2 = \frac{2n(n + 1)(2n + 1)}{3}$ and so that 

 
            \begin{enumerate}[label=(\roman*)]
                \item $1^2 + 3^2 + \ldots +(2n + 1)^2 = (1^2 + 2^2 + \ldots +(2n + 1)^2) - (2^2 + 4^2 + \ldots + (2n)^2)$\item $= \frac{(2n + 1)(2n + 2)(4n + 3)}{6} - \frac{2n(n + 1)(2n + 1)}{3} = \frac{(n + 1)(2n + 1)(2n + 3)}{3}$.  Thus,
            \end{enumerate}
            $|a - b| = \left| 2\cdot \frac{44\cdot23\cdot45}{3} - 2\cdot \frac{22 \cdot 43 \cdot 45}{3} - 45^2 + 20\right| = |-25| =\boxed{025}$.

 





\textbf{Solution 4}

Let $\Delta n$ denote the sum $1+2+3+ \dots +n-1+n$. We can easily see from the fact "It is well-known that $\tau(n)$ is odd if and only if $n$ is a perfect square.", that 

 $a = (2^2-1^2) + (4^2-3^2) \dots (44^2 - 43^2) = (2+1)(2-1)+(4+3)(4-3) \dots (44+43)(44-43) = 1+2+3...44 = \Delta 44$. 

 $b = 3^2-2^2+5^2-4^2...2006-44^2 = (3+2)(3-2) \dots (43+42)(43-42) + 1 + (2005 - 44^2) = \Delta 43 + 69$. 

 $a-b = \Delta 44-\Delta 43-69 = 44-69 = -25$. They ask for $|a-b|$, so our answer is $|-25| = \boxed{025}$

 -Alexlikemath

 




	\item (\textit{AIME I 2005}) A particle moves in the Cartesian Plane according to the following rules:


 
            \begin{enumerate}
                \item From any lattice point $(a,b),$ the particle may only move to $(a+1,b), (a,b+1),$ or $(a+1,b+1).$
\item There are no right angle turns in the particle's path.
            \end{enumerate}
            How many different paths can the particle take from $(0,0)$ to $(5,5)$?


 
 \textbf{Answer: }083



\textbf{Solution 1}

The length of the path (the number of times the particle moves) can range from $l = 5$ to $9$; notice that $d = 10-l$ gives the number of diagonals. Let $R$ represent a move to the right, $U$ represent a move upwards, and $D$ to be a move that is diagonal. \href{https://artofproblemsolving.com/wiki/index.php?title=Casework}{Casework} upon the number of diagonal moves:

 
                \begin{itemize}
                    \item \textbf{Case } $d = 1$: It is easy to see only $2$ cases.\item \textbf{Case } $d = 2$: There are two diagonals. We need to generate a string with $3$ $R$'s, $3$ $U$'s, and $2$ $D$'s such that no two $R$'s or $U$'s are adjacent. The $D$'s split the string into three sections ($-D-D-$): by the \href{https://artofproblemsolving.com/wiki/index.php?title=Pigeonhole\_principle}{Pigeonhole principle} all of at least one of the two letters must be all together (i.e., stay in a row).
                \end{itemize}


                
            \begin{enumerate}[label=(\roman*)]
                \item If both $R$ and $U$ stay together, then there are $3 \cdot 2=6$ ways.\item If either $R$ or $U$ splits, then there are $3$ places to put the letter that splits, which has $2$ possibilities. The remaining letter must divide into $2$ in one section and $1$ in the next, giving $2$ ways. This totals $6 + 3\cdot 2\cdot 2 = 18$ ways.
            \end{enumerate}
            
                \begin{itemize}
                    \item \textbf{Case } $d = 3$: Now $2$ $R$'s, $2$ $U$'s, and $3$ $D$'s, so the string is divided into $4$ partitions ($-D-D-D-$).
                \end{itemize}


                
            \begin{enumerate}[label=(\roman*)]
                \item If the $R$'s and $U$'s stay together, then there are $4 \cdot 3 = 12$ places to put them.\item If one of them splits and the other stays together, then there are $4 \cdot {3\choose 2}$ places to put them, and $2$ ways to pick which splits, giving $4 \cdot 3 \cdot 2 = 24$ ways.\item If both groups split, then there are ${4\choose 2}=6$ ways to arrange them. These add up to $12 + 24 + 6 = 42$ ways.
            \end{enumerate}
            
                \begin{itemize}
                    \item \textbf{Case } $d = 4$: Now $1$ $R$, $1$ $U$, $4$ $D$'s ($-D-D-D-D-$). There are $5$ places to put $R$, $4$ places to put $U$, giving $20$ ways.\item \textbf{Case } $d = 5$: It is easy to see only $1$ case.
                \end{itemize}


                Together, these add up to $2 + 18 + 42 + 20 + 1 = \boxed{083}$.

 





\textbf{Solution 2}

Another possibility is to use block-walking and \href{https://artofproblemsolving.com/wiki/index.php?title=Recursion}{recursion}: for each vertex, the number of ways to reach it is $a + b + c$, where $a$ is the number of ways to reach the vertex from the left (without having come to \textit{that} vertex (the one on the left) from below), $b$ is the number of ways to reach the vertex from the vertex diagonally down and left, and $c$ is the number of ways to reach the vertex from below (without having come to \textit{that} vertex (the one below) from the left).

 Assign to each point $(i,j)$ the triplet $(a_{i,j}, b_{i,j}, c_{i,j})$. Let $s(i,j) = a_{i,j}+ b_{i,j}+ c_{i,j}$. Let all lattice points that contain exactly one negative coordinate be assigned to $(0,0,0)$. This leaves the lattice points of the first quadrant, the positive parts of the $x$ and $y$ axes, and the origin unassigned. As a seed, assign to $(0,1,0)$. (We will see how this correlates with the problem.) Then define for each lattice point $(i,j)$ its triplet thus:
 \begin{align*} a_{i,j} &= s(i-1,j) - c_{i-1,j}\\ b_{i,j} &= s(i-1,j-1) \\ c_{i,j} &= s(i,j-1) - a_{i,j-1}. \end{align*}
It is evident that $s(i,j)$ is the number of ways to reach $(i,j)$ from $(0,0)$. Therefore we compute vertex by vertex the triplets $(a_{i,j}, b_{i,j}, c_{i,j})$ with $0 \leq i, j \leq 5$. 
                    \begin{figure}[H]
                        \centering
                        \includegraphics[width=0.5\textwidth]{Solution/AIME_I_2005_13_Solution_Figure_1.png}
                    \end{figure}
                    Finally, after simple but tedious calculations, we find that $(a_{5,5}, b_{5,5}, c_{5,5}) = (28,27,28)$, so $s(i,j)=28+27+28 = \boxed{083}$.

 




	\item (\textit{AIME I 2005}) Consider the points $A(0,12), B(10,9), C(8,0),$ and $D(-4,7).$ There is a unique square $S$ such that each of the four points is on a different side of $S.$ Let $K$ be the area of $S.$ Find the remainder when $10K$ is divided by 1000.


 
 \textbf{Answer: }936



\textbf{Solution}


                    \begin{figure}[H]
                        \centering
                        \includegraphics[width=0.5\textwidth]{Solution/AIME_I_2005_14_Solution_Figure_1.png}
                    \end{figure}
                    





\textbf{Solution 1}

Consider a point $E$ such that $AE$ is \href{https://artofproblemsolving.com/wiki/index.php?title=Perpendicular}{perpendicular} to $BD$, $AE$ intersects $BD$, and $AE = BD$. E will be on the same side of the square as point $C$.

 Let the coordinates of $E$ be $(x_E,y_E)$. Since $AE$ is perpendicular to $BD$, and $AE = BD$, we have $9 - 7 = x_E - 0$ and $10 - ( - 4) = 12 - y_E$The coordinates of $E$ are thus $(2, - 2)$.

 Now, since $E$ and $C$ are on the same side, we find the slope of the sides going through $E$ and $C$ to be $\frac { - 2 - 0}{2 - 8} = \frac {1}{3}$. Because the other two sides are perpendicular, the slope of the sides going through $B$ and $D$ are now $- 3$.

 Let $A_1,B_1,C_1,D_1$ be the vertices of the square so that $A_1B_1$ contains point $A$, $B_1C_1$ contains point $B$, and etc. Since we know the slopes and a point on the line for each side of the square, we use the point slope formula to find the linear equations. Next, we use the equations to find $2$ vertices of the square, then apply the distance formula. 

 We find the coordinates of $C_1$ to be $(12.5,1.5)$ and the coordinates of $D_1$ to be $( - 0.7, - 2.9)$. Applying the distance formula, the side length of our square is $\sqrt {\left( \frac {44}{10} \right)^2 + \left( \frac {132}{10} \right)^2} = \frac {44}{\sqrt {10}}$.

 Hence, the area of the square is $K = \frac {44^2}{10}$. The remainder when $10K$ is divided by $1000$ is $936$.

 





\textbf{Solution 2}

Let $(a,b)$ denote a \href{https://artofproblemsolving.com/wiki/index.php?title=Normal\_vector&action=edit&redlink=1}{normal vector} of the side containing $A$. Note that $\overline{AC}, \overline{BD}$ intersect and hence must be opposite \href{https://artofproblemsolving.com/wiki/index.php?title=Vertex}{vertices} of the square. The lines containing the sides of the square have the form $ax+by=12b$, $ax+by=8a$, $bx-ay=10b-9a$, and $bx-ay=-4b-7a$. The lines form a square, so the distance between $C$ and the line through $A$ equals the distance between $D$ and the line through $B$, hence $8a+0b-12b=-4b-7a-10b+9a$, or $-3a=b$. We can take $a=-1$ and $b=3$. So the side of the square is $\frac{44}{\sqrt{10}}$, the area is $K=\frac{1936}{10}$, and the answer to the problem is $\boxed{936}$.

 





\textbf{Solution 3}

Let $\begin{bmatrix} a \\ b \end{bmatrix}$ be the unit vector parallel to one side of the square. The unit vector parallel to the perpendicular side is then  $\begin{bmatrix} b \\ -a \end{bmatrix}$. The dot product of a diagonal of quadrilateral $ABCD$ with the corresponding unit vector gives the side length. Since the side lengths of the square are equal, $\overrightarrow{DB}\cdot \begin{bmatrix} a \\ b \end{bmatrix}= \overrightarrow{AC}\cdot \begin{bmatrix} b \\ -a \end{bmatrix}$. Plugging in $\overrightarrow{AC}=\begin{bmatrix} 8 \\ -12 \end{bmatrix}$ and $\overrightarrow{DB}=\begin{bmatrix} 14 \\ 2 \end{bmatrix}$ yields $a=3b$. Since $\begin{bmatrix} a \\ b \end{bmatrix}$ is a unit vector, it is equal to $\begin{bmatrix} \frac{3}{\sqrt{10}} \\ \frac{1}{\sqrt{10}} \end{bmatrix}$. Taking the dot product with $\overrightarrow{DB}$ gives the side length, which is $\frac {44}{\sqrt {10}}$, so the area is $\frac {44^2}{10}$, and the answer is $\boxed{936}$.

 




	\item (\textit{AIME I 2005}) Triangle $ABC$ has $BC=20.$ The incircle of the triangle evenly trisects the median $AD.$ If the area of the triangle is $m \sqrt{n}$ where $m$ and $n$ are integers and $n$ is not divisible by the square of a prime, find $m+n.$


 
 \textbf{Answer: }038



\textbf{Solution 1}


                    \begin{figure}[H]
                        \centering
                        \includegraphics[width=0.5\textwidth]{Solution/AIME_I_2005_15_Solution_Figure_1.png}
                    \end{figure}
                    Let $E$, $F$ and $G$ be the points of tangency of the incircle with $BC$, $AC$ and $AB$, respectively.  Without loss of generality, let $AC < AB$, so that $E$ is between $D$ and $C$.  Let the length of the median be $3m$.  Then by two applications of the \href{https://artofproblemsolving.com/wiki/index.php?title=Power\_of\_a\_Point\_Theorem}{Power of a Point Theorem}, $DE^2 = 2m \cdot m = AF^2$, so $DE = AF$.  Now, $CE$ and $CF$ are two tangents to a circle from the same point, so by the Two Tangent Theorem $CE = CF = c$ and thus $AC = AF + CF = DE + CE = CD = 10$.  Then $DE = AF = AG = 10 - c$ so $BG = BE = BD + DE = 20 - c$ and thus $AB = AG + BG = 30 - 2c$.

 Now, by \href{https://artofproblemsolving.com/wiki/index.php?title=Stewart\%27s\_Theorem}{Stewart's Theorem} in triangle $\triangle ABC$ with \href{https://artofproblemsolving.com/wiki/index.php?title=Cevian}{cevian} $\overline{AD}$, we have

 
 \[(3m)^2\cdot 20 + 20\cdot10\cdot10 = 10^2\cdot10 + (30 - 2c)^2\cdot 10.\]

 Our earlier result from Power of a Point was that $2m^2 = (10 - c)^2$, so we combine these two results to solve for $c$ and we get

 
 \[9(10 - c)^2 + 200 = 100 + (30 - 2c)^2 \quad \Longrightarrow \quad c^2 - 12c + 20 = 0.\]

 Thus $c = 2$ or $= 10$.  We discard the value $c = 10$ as extraneous (it gives us a line) and are left with $c = 2$, so our triangle has area $\sqrt{28 \cdot 18 \cdot 8 \cdot 2} = 24\sqrt{14}$ and so the answer is $24 + 14 = \boxed{038}$.

 





\textbf{Solution 2}

WLOG let E be be between C & D (as in solution 1). Assume $AD = 3m$. We use power of a point to get that $AG = DE = \sqrt{2}m$ and  $AB = AG + GB = AG + BE = 10+2\sqrt{2} m$ 

 Since now we have $AC = 10$, $BC = 20, AB = 10+2\sqrt{2} m$ in triangle $\triangle ABC$ and \href{https://artofproblemsolving.com/wiki/index.php?title=Cevian}{cevian} $AD = 3m$.  Now, we can apply \href{https://artofproblemsolving.com/wiki/index.php?title=Stewart\%27s\_Theorem}{Stewart's Theorem}.

 
 \[2000 + 180 m^2 = 10(10+2\sqrt{2}m)^{2} + 1000\]
 \[1000 + 180 m^2 = 1000 + 400\sqrt{2}m + 80 m^{2}\]
 \[100 m^2 = 400\sqrt{2}m\]

 $m = 4\sqrt{2}$ or $m = 0$ if $m = 0$, we get a degenerate triangle, so $m = 4\sqrt{2}$, and thus $AB = 26$. You can now use \href{https://artofproblemsolving.com/wiki/index.php?title=Heron\%27s\_Formula}{Heron's Formula} to finish. The answer is $24 \sqrt{14}$, or $\boxed{038}$.

 -Alexlikemath

 





\textbf{Solution 3}

Let $E, F$, and $G$ be the point of tangency (as stated in Solution 1). We can now let $AD$ be $3m$. By using \href{https://artofproblemsolving.com/wiki/index.php?title=Power\_of\_a\_Point\_Theorem}{Power of a Point Theorem} on A to the incircle, you get that $AG^2 = 2m^2$. We can use it again on point D to the incircle to get the equation $(10 - CE)^2 = 2m^2$. Setting the two equations equal to each other gives $(10 - CE)^2 = AG^2$, and it can be further simplified to be $10 - CE = AG$

 Let lengths $AC$ and $AB$ be called $b$ and $c$, respectively. We can write $AG$ as $\frac{b + c - 20}{2}$ and $CE$ as $\frac{b + 20 - c}{2}$. Plugging these into the equation, you get:

 
 \[10 - \frac{b + 20 - c}{2} = \frac{b + c - 20}{2}\]
 \[b + c - 20 + b + 20 - c = 20 \rightarrow b = 10\]

 Additionally, by \href{https://artofproblemsolving.com/wiki/index.php?title=Median\_of\_a\_triangle}{Median of a triangle} formula, you get that $3m = \frac{\sqrt{2c^2 - 200}}{2}$

 Refer back to the fact that $AG^2 = 2m^2$. We can now plug in our variables.

 
 \[\left(\frac{b + c - 20}{2}\right)^2 = 2m^2 \rightarrow (c - 10)^2 = 8m^2\]
 \[c^2 - 20c + 100 = 8 \cdot \frac{2c^2 - 200}{36}\]
 \[9c^2 - 180c + 900 = 4c^2 - 400\]
 \[5c^2 - 180c + 1300 = 0\]
 \[c^2 - 36c + 260 = 0\]

 Solving, you get that $c = 26$ or $10$, but the latter will result in a degenerate triangle, so $c = 26$.Finally, you can use \href{https://artofproblemsolving.com/wiki/index.php?title=Heron\%27s\_Formula}{Heron's Formula} to get that the area is $24\sqrt{14}$, giving an answer of $\boxed{038}$

 




\end{enumerate}